\section{Heuristic Agent}

\subsection{Current Implementations}

The repository contains submission adapters for \texttt{EnterpriseHeuristicAgentV11a} and \texttt{EnterpriseHeuristicAgentV11b}. The v11a adapter describes the agent as a rule-based priority heuristic with Restore-only threat response and preemptive Operational Zone blocking. The v11b adapter adds a redesigned messaging protocol for coordinated blocking.

Because local documents report mixed results, this draft treats v11a as the best documented baseline and v11b as an experimental messaging variant until a fresh benchmark confirms otherwise.

\subsection{Design Philosophy}

The heuristic is deterministic. At every step, it evaluates a priority cascade and executes the first valid action. There is no learned value function and no stochastic exploration. The strength of the agent comes from encoding environment-specific timing facts:

\begin{itemize}
  \item \textbf{Restore-biased response}: \texttt{Remove} now clears high-density unsigned malware files in this fork, but it still cannot evict privileged/root red sessions. \texttt{Restore} remains the reliable eviction path after privilege escalation.
  \item \textbf{Decoy saturation}: each host is filled with up to three decoys, increasing the chance that blind red exploitation wastes time.
  \item \textbf{Malicious-file detection}: high-density unsigned files are treated as strong evidence of compromise.
  \item \textbf{Flag-age filtering}: transient process flags are delayed to avoid green false positives.
  \item \textbf{Phase-aware blocking}: operational-zone paths are blocked around phase transitions to avoid windows where red can move while blue is busy.
\end{itemize}

\subsection{v11a: Preemptive Operational-Zone Blocking}

The v11a source states that the relevant transition windows are:

\begin{itemize}
  \item before Phase 1, block \texttt{restricted\_zone\_a\_subnet -> operational\_zone\_a\_subnet};
  \item before Phase 2, block \texttt{restricted\_zone\_b\_subnet -> operational\_zone\_b\_subnet}.
\end{itemize}

The rationale is timing. If the agent waits until a phase transition is visible, it may be occupied by higher-priority Restore actions. Preemptive blocking pays a small service-availability cost before the phase but avoids a high-risk open path at the beginning of the active mission phase.

\subsection{v11b: Messaging Redesign}

The v11b source describes a new 8-bit message format with fields for red detection, root-level suspicion, restore-busy status, zone-clear status, threat count, and block requests. Its goal is to let upstream agents ask downstream operational-zone agents to block traffic from compromised zones.

The research notes warn that earlier messaging designs showed little measurable value. The likely reasons are:
\begin{itemize}
  \item \texttt{comms\_policy} already provides much of the useful coordination;
  \item only a small subset of phase-agent combinations has useful upstream dependency;
  \item one-step message delay is large relative to red action timing.
\end{itemize}

\subsection{Submission Adapter}

The heuristic submission wraps the inner agent with a \texttt{BlueFlatWrapperV2}-based environment. The adapter bridges the official \texttt{evaluation.py} interface by:

\begin{itemize}
  \item retrieving action masks from the wrapper;
  \item forwarding the action index returned by the heuristic;
  \item storing outgoing messages from each agent;
  \item injecting those messages into the next environment step, because the official evaluator does not pass them directly.
\end{itemize}
